\section{Toy endpoint--boundary operator}
\label{sec:toy-boundary-operator}
% Status: conjectural/open. Translates mechanisms into named matrix entries and derivation obligations.

The mechanism sections reduce to a common two-channel form.  This section gives
the minimal endpoint--boundary operator used to name the matrix entries.

Let \(h_J(x)\) denote a boundary-localized order-parameter amplitude and \(a_J(x)\) denote an adjoint/current amplitude.  The quadratic light-sector action is
\begin{equation}
  S^{(2)}_J
  =
  \int d^4x\,
  \begin{pmatrix}h_J^\dagger & a_J^\dagger\end{pmatrix}
  \begin{pmatrix}
    \lambda & -\kappa_J\\
    -\kappa_J & \lambda+\tau_J
  \end{pmatrix}
  \begin{pmatrix}h_J\\ a_J\end{pmatrix}.
  \label{eq:toy-boundary-action}
\end{equation}
Here \(\lambda\) is the spectral variable after normalization.  The DeVries target is
\begin{equation}
  \tau_J=J,\qquad \kappa_J^2=J,
  \label{eq:toy-devries-target}
\end{equation}
which gives
\begin{equation}
  \det
  \begin{pmatrix}
    \lambda & -\sqrt J\\
    -\sqrt J & \lambda+J
  \end{pmatrix}
  =
  \lambda^2+J\lambda-J.
  \label{eq:toy-devries-det}
\end{equation}

The fields \(h_J\) and \(a_J\) are placeholders with fixed jobs.  The \(h_J\) channel must carry the doublet/order-parameter invariant when \(J=3/4\).  The \(a_J\) channel must carry the adjoint/current invariant when \(J=2\).  The off-diagonal entry \(\kappa_J\) is the physical mixing amplitude between the two channels.  The diagonal entry \(\tau_J\) is the trace datum of the adjoint/current side.

\subsection{Boundary origin}

Interval field theory supplies a mechanism for such a matrix.  Boundary mass and kinetic terms modify the variational boundary conditions, and brane-localized kinetic terms can make the boundary condition depend on the four-dimensional eigenvalue \cite{CsakiHubiszMeade2005}.  A schematic bulk-plus-boundary system is
\begin{equation}
  \begin{aligned}
  S_{\rm bulk}
  =
  &\int d^4x\int_0^Ldy\,
  \frac12\Big[(\partial_M A_J)^2+(\partial_M H_J)^2\Big]\\
  &+S_{\rm bdry}[A_J,H_J] .
  \end{aligned}
  \label{eq:bulk-boundary-action}
\end{equation}
After solving the bulk equations and restricting to the two light boundary amplitudes, the effective inverse propagator has the form
\begin{equation}
  K^{\rm eff}_J(\lambda)=
  \begin{pmatrix}
    \lambda+\Sigma_{hh,J}(\lambda) & \Sigma_{ha,J}(\lambda)\\
    \Sigma_{ah,J}(\lambda) & \lambda+\Sigma_{aa,J}(\lambda)
  \end{pmatrix}.
  \label{eq:effective-boundary-kernel}
\end{equation}
The DeVries calculation asks for a normalization and a boundary limit in which
\begin{equation}
  \Sigma_{hh,J}=0,\qquad
  \Sigma_{aa,J}=J,\qquad
  \Sigma_{ha,J}\Sigma_{ah,J}=J.
  \label{eq:self-energy-target}
\end{equation}

This formulation gives each matrix entry a physical address.  \(\Sigma_{hh,J}\) tests whether the order-parameter channel can be used as the zero reference.  \(\Sigma_{aa,J}\) tests whether the adjoint/current channel supplies the trace \(J\).  \(\Sigma_{ha,J}\) tests whether endpoint, brane, or overlap physics supplies \(\sqrt J\).

\subsection{Endpoint interpretation}

In an open-string or D-brane reading, \(a_J\) is an endpoint current or Chan--Paton-like gauge excitation, while \(h_J\) is a boundary order-parameter mode.  Coincident branes naturally give nonabelian gauge fields from endpoint sectors \cite{Polchinski1994,TongString}.  The toy model demands a more specific result: the endpoint algebra must reduce to the entries in Eq.~\eqref{eq:toy-devries-target}.

The corresponding route is
\begin{equation}
\begin{array}{ccc}
\text{endpoint labels} & \longrightarrow & \kappa_J^2=J\\
\text{adjoint/current trace} & \longrightarrow & \tau_J=J\\
\text{boundary normalization} & \longrightarrow & \Sigma_{hh,J}=0 .
\end{array}
\label{eq:endpoint-route-table}
\end{equation}
This table names the calculation required of an endpoint or boundary source.

\subsection{Electroweak matching}

The toy operator must match the electroweak projective data.  Boundary data
must reduce to \(\mathcal M_J\); diagonalizing \(\mathcal M_J\) then supplies
\(\{x_+(J),x_-(J)\}\), the W/Z pole-ratio reading, and the partner-branch
condition.
The positive branch should yield the pole-scheme target
\begin{equation}
  \frac{M_{W,\rm pole}^2}{M_{Z,\rm pole}^2}=
  \frac{x_+(J_H)}{x_+(J_{\rm adj})}
  +\Delta_{\rm toy},
  \qquad
  (J_H,J_{\rm adj})=(3/4,2).
  \label{eq:toy-ew-match}
\end{equation}
Here \(\Delta_{\rm toy}\) records every step missing from the toy model:
field normalization, source scale, effective matching, and pole-scheme
placement.
The negative branch should identify the scalar or boundary partner in the same
basis.  The toy operator therefore links ordered electroweak sampling, the
determinant calculation, and the scalar-branch interpretation.
